% TOP_PROBLEM: {"problemId": "problem.yang-mills-existence-and-mass-gap", "canonicalTitle": "Yang-Mills Existence and Mass Gap", "releaseRank": 3, "primaryDomain": "mathematical_physics", "primaryDomainLabel": "Mathematical physics", "displayStatus": "open", "releaseStatus": "open"}
% !TeX program = lualatex
% Definition and sources only; this file is not an LLM solution attempt.
\documentclass[11pt]{article}
\usepackage[margin=25mm]{geometry}
\usepackage{fontspec}
\setmainfont{Latin Modern Roman}
\usepackage{amsmath,amssymb,mathtools}
\usepackage{unicode-math}
\setmathfont{Latin Modern Math}
\usepackage{xurl,hyperref}
\hypersetup{colorlinks=true,urlcolor=blue}
\setcounter{secnumdepth}{0}
\setlength{\parindent}{0pt}
\setlength{\parskip}{5pt}
\setlength{\emergencystretch}{3em}
\begin{document}
\section*{\texorpdfstring{Yang-Mills Existence and Mass Gap}{Yang-Mills Existence and Mass Gap}}\label{3}
\subsection{Definitions and mathematical statement}
A classical $G$-connection has curvature $F_A=dA+A\wedge A$ and Yang--Mills action proportional to $\int_{{\mathbb{R}}^4}\langle F_A,F_A\rangle\,dx$. The question concerns a rigorous, nontrivial \emph{quantum} theory associated with this action, not only classical solutions or a formal functional integral. In its vacuum representation, let $\Omega$ be the vacuum and $H\ge0$ the self-adjoint generator of time translations, with $H\Omega=0$. A positive finite mass gap means
\[
 0<m:=\inf\bigl(\operatorname{spec}(H)\setminus\{0\}\bigr)<\infty.
\]
For every compact simple $G$, construct such a theory with the local observable fields, positivity, covariance, locality, and short-distance conditions required in Jaffe--Witten, Sections 3--4. Those axioms are part of the target and are not replaced by the displayed spectral condition alone.
\par\smallskip\textit{Formulation and concept references:} \href{https://www.claymath.org/wp-content/uploads/2022/06/yangmills.pdf}{[S1]}.

\subsection{Short English statement}
Can every compact simple gauge symmetry produce a rigorously defined four-dimensional Yang--Mills quantum theory with a nonzero energy gap above its vacuum?
\subsection{Sources}
\begin{itemize}
\item[S1]Arthur Jaffe and Edward Witten, Quantum Yang-Mills Theory, official Clay problem description.
\url{https://www.claymath.org/wp-content/uploads/2022/06/yangmills.pdf}.
\textit{Formulation source. Consulted 24 September 2026: Sections 3--4, especially page 6.}
\item[S2]Yang-Mills \& the Mass Gap.
\url{https://www.claymath.org/millennium/yang-mills-the-maths-gap/}.
\textit{Further reference; not a proof endorsement.}
\item[S3]Towards a Proof of Mass Gap in 3d Yang-Mills Theory — V. P. Nair.
\url{https://arxiv.org/abs/2608.10133}.
\textit{Further reference; not a proof endorsement.}
\item[S4]A Constructive Proof of Existence and Mass Gap for Pure SU(3) Yang-Mills in Four-Dimensional Space-Time — D. C. Jacobsen.
\url{https://arxiv.org/abs/2506.00284}.
\textit{Further reference; not a proof endorsement.}
\end{itemize}
\end{document}
